% Generated draft: requires mathematical and source review.
\documentclass[11pt]{article}
\usepackage[margin=25mm]{geometry}
\usepackage{fontspec}
\setmainfont{Latin Modern Roman}
\usepackage{amsmath,amssymb,mathtools,unicode-math}
\setmathfont{Latin Modern Math}
\usepackage{enumitem,needspace,xurl,hyperref}
\setlength{\parindent}{0pt}
\setlength{\parskip}{5pt}
\setlength{\emergencystretch}{3em}
\newcommand{\N}{\mathbb{N}}
\newcommand{\Z}{\mathbb{Z}}
\newcommand{\Q}{\mathbb{Q}}
\newcommand{\R}{\mathbb{R}}
\newcommand{\C}{\mathbb{C}}
\newcommand{\F}{\mathbb{F}}
\newcommand{\A}{\mathbb{A}}
\newcommand{\PP}{\mathbb{P}}
\newcommand{\E}{\mathbb{E}}
\newcommand{\eps}{\varepsilon}
\newcommand{\dd}{\,\mathrm{d}}
\begin{document}
\section*{UnsolvedMath ID: 30002553}
\textbf{Source problem code:} catalogue entry
\subsection{Definitions and mathematical statement}
The supplied record concerns rotationally symmetric manifolds \(M^n\) and \(N^\ell\), with metrics written as
\[
\dd r^2+h(r)^2\dd\omega^2,
\qquad
\dd\phi^2+g(\phi)^2\dd\chi^2,
\]
where \(n\geq 3\), \(k\geq 1\), and \(\ell=k(k+n-2)\). Here \(r\) and \(\phi\) denote the respective radial coordinates, and the angular metric factors are denoted by \(\dd\omega^2\) and \(\dd\chi^2\). The same symbol \(\phi\) denotes the scalar profile of an equivariant wave map. Primes denote differentiation of \(h\) and \(g\) with respect to their arguments.

For a source manifold \(M\) that is \emph{admissible in the sense of the cited report}, the displayed profile equation is
\[
\partial_t^2\phi-\partial_r^2\phi
-(n-1)\frac{h'(r)}{h(r)}\partial_r\phi
+\frac{\ell}{h(r)^2}g(\phi)g'(\phi)=0,
\]
with Cauchy data
\[
\phi(0,r)=\phi_0(r),\qquad
\partial_t\phi(0,r)=\phi_1(r).
\]
The record designates \(H^{n/2}\times H^{n/2-1}\) as the critical data space.

\textbf{Question (unconditional uniqueness).} For the admissible geometries and equivariance classes intended by the report, do any two solutions of this Cauchy problem with the same critical Sobolev data coincide throughout their common interval of existence, assuming only the intended critical Sobolev solution regularity and imposing no additional Strichartz-space condition?

More explicitly, if \(I\) is a common time interval containing \(0\), and \(\phi\) and \(\psi\) are solutions in the intended critical solution class with
\[
\bigl(\phi,\partial_t\phi\bigr)(0)
=
\bigl(\psi,\partial_t\psi\bigr)(0)
=(\phi_0,\phi_1)
\in H^{n/2}\times H^{n/2-1},
\]
is it necessarily the case that \(\phi=\psi\) on their common spacetime domain, without separately assuming membership in a Strichartz space?

\textbf{Unresolved qualifications.} This is a conditional formulation of the supplied question. The record does not define admissibility, the radial coordinate domains, the hypotheses on \(h\) and \(g\), the angular geometry or equivariant ansatz, or the required pole, endpoint, and spatial-infinity conditions. It also does not specify the underlying measure and geometric realization of the Sobolev spaces, the time regularity of the solution class, or the sense in which the equation is satisfied. Consequently, the phrase ``intended critical solution class'' cannot be replaced by a fully specified function space from this record alone. The notation \(N^\ell\) and the simultaneous use of \(\ell=k(k+n-2)\) as the equation coefficient are retained exactly as supplied; the relationship between target dimension and angular spectral data requires clarification. These missing qualifications prevent a fully determined uniqueness conjecture.
\subsection{Short English statement}
The supplied question asks whether two equivariant wave-map solutions with the same data in \(H^{n/2}\times H^{n/2-1}\) must agree using critical Sobolev regularity alone, without an extra Strichartz-space assumption. The admissible geometries and precise solution class are not defined in the supplied record, so the statement remains conditional on those missing specifications.
\subsection{Sources}
\begin{itemize}
\item Supplied local UnsolvedMath registry record: permanent ID \(30002553\), source problem code \texttt{catalogue entry}, title \emph{Unconditional Uniqueness for Critical Equivariant Wave Maps}, in the Oberwolfach Reports Open Problems collection. No individual proposer is identified; the supplied proposed year is 2014.
\item \emph{Nonlinear Evolution Equations: Analysis and Numerics} (2014), Oberwolfach Reports, pp.\ 781--868, DOI \texttt{10.4171/owr/2014/14}. Supplied source and literature reference: \url{https://doi.org/10.4171/owr/2014/14}. The background describes these reports as records of workshops and their open-problem discussions. No precise problem page or contribution author is supplied.
\item D'Ancona--Zhang, IMRN, supplied literature reference: \url{https://doi.org/10.1093/imrn/rnv152}. The record does not supply an article title, theorem number, or a complete statement of the relevant uniqueness result.
\item Literature assessment dated exactly \texttt{2026-08-21}: the supplied curation labels the problem open and states that unconditional uniqueness at the stated critical regularity is not established in the rotationally symmetric setting, while D'Ancona--Zhang obtain uniqueness after increasing regularity. This is supplied dated curation, not an independently verified literature assessment. The background expressly requests independent verification of time-sensitive status and cited claims.
\item The registry labels its statement \texttt{corrected\_verified} and reports that it was checked against the official Oberwolfach report to restore hypotheses and notation. This is a supplied verification claim; no independent source inspection or mathematical verification is asserted here.
\end{itemize}
\end{document}
